\section{Experimental Evaluation}
\label{sec:experiments}

We evaluate our proposed continuous wave-theoretic merging framework, \textbf{PhaseMerge}, on vision transformers across highly conflicting datasets. Below, we describe our experimental setup, present our key findings, and analyze the resulting metrics under various deployment configurations.

\subsection{Experimental Setup}

\textbf{Backbone Model and Layer Selection.} We use the Vision Transformer Tiny architecture (\texttt{vit\_tiny\_patch16\_224}) \cite{dosovitskiy2020image} containing $L = 12$ transformer blocks as our backbone. Following standard practice in weight-space merging, we target the dense weight matrices in the Multi-Head Self-Attention (`qkv`, `proj`) and Multi-Layer Perceptron (`fc1`, `fc2`) modules, yielding a total of 48 targeted dense weight matrices.

\textbf{Downstream Tasks.} We evaluate on four highly conflicting vision tasks: MNIST, FashionMNIST, CIFAR-10, and SVHN. To establish highly specialized task experts, each model is fine-tuned individually on its target task (500 training samples) and evaluated on all four task test sets (100 samples per task) to verify extreme task conflict and ensure statistically robust accuracy reporting. Individual expert diagonal performance across 3 seeds reaches $81.00 \pm 0.82\%$ on MNIST, $74.67 \pm 1.25\%$ on FashionMNIST, $71.67 \pm 4.03\%$ on CIFAR-10, and $85.33 \pm 2.87\%$ on SVHN, while off-diagonal performance collapses to random guessing (averaging $\approx 10\%$). This extreme off-diagonal failure profile highlights the highly conflicting nature of our multi-task setting, presenting a severe challenge for naive parameter merging.

\textbf{Optimization and Hyperparameters.} For PhaseMerge, we optimize a $2 \times 2$ continuous phase grid per task and layer ($r = 2$, 768 parameters) or a uniform phase rotation per task and layer ($r = 1$, 192 parameters). Optimization is carried out via the Adam optimizer with a learning rate of $1\times 10^{-2}$ for 5 steps to ensure ultra-fast test-time convergence. Unsupervised calibration is driven by minimizing prediction entropy over small calibration streams $D_{\text{cal}}$ of sizes $M \in \{4, 16, 32\}$. All metrics represent the mean and standard deviation across 3 independent random seeds.

\textbf{Baselines.} We compare PhaseMerge against several established baselines:
\begin{itemize}
    \item \textbf{Uniform Task Arithmetic (TA) \cite{ilharco2022editing}:} Linear model merging with a static weight coefficient of $0.3$.
    \item \textbf{FREE-Merging \cite{freemerge2024}:} A static, non-adaptive frequency-domain model-merging baseline that applies a hard-coded low-pass Fourier filter (retaining the lowest $85\%$ of frequencies) to task vectors prior to linear merging.
    \item \textbf{AdaMerging \cite{yang2023adamerging}:} Unconstrained adaptive layer-wise parameter optimization in real space (48-dimensional).
    \item \textbf{PolyMerge \cite{polymerge2024}:} A highly constrained adaptive merging baseline that parameterizes layer coefficients using a quadratic depth-wise polynomial (12-dimensional).
\end{itemize}

\subsection{Multi-Task Merging Accuracy}

We first report the multi-task average accuracy across the test sets of all four tasks. All optimization-based methods are calibrated using $M=16$ samples. We evaluate performance under three distinct schema regimes: unquantized FP32, 8-bit Post-Training Quantization (PTQ), and 4-bit PTQ.

\begin{table}[ht]
\caption{Multi-Task Average Accuracy (\%) across MNIST, FashionMNIST, CIFAR-10, and SVHN test sets under various quantization regimes. Boldface denotes the best overall performance.}
\label{tab:main_results}
\vskip 0.15in
\begin{center}
\begin{small}
\begin{sc}
\resizebox{\columnwidth}{!}{%
\begin{tabular}{lccc}
\toprule
Merging Method & FP32 & 8-bit PTQ & 4-bit PTQ \\
\midrule
Uniform TA & 38.25 $\pm$ 1.34\% & 37.75 $\pm$ 1.43\% & 33.17 $\pm$ 1.39\% \\
FREE-Merging & 27.17 $\pm$ 1.96\% & 27.17 $\pm$ 2.18\% & 24.33 $\pm$ 2.37\% \\
AdaMerging & 42.00 $\pm$ 0.89\% & 41.67 $\pm$ 1.45\% & 37.50 $\pm$ 1.22\% \\
PolyMerge & \textbf{48.00 $\pm$ 1.62\%} & \textbf{48.00 $\pm$ 1.47\%} & \textbf{43.42 $\pm$ 1.30\%} \\
U-PhaseMerge ($r=1$, Ours) & 42.83 $\pm$ 1.76\% & 42.33 $\pm$ 1.76\% & 37.42 $\pm$ 1.94\% \\
PhaseMerge ($r=2$, Ours) & 40.75 $\pm$ 1.43\% & 40.83 $\pm$ 1.18\% & 36.92 $\pm$ 0.92\% \\
\bottomrule
\end{tabular}%
}
\end{sc}
\end{small}
\end{center}
\vskip -0.1in
\end{table}

As shown in Table~\ref{tab:main_results}, PolyMerge \cite{polymerge2024} achieves the highest overall multi-task average accuracy ($48.00 \pm 1.62\%$ FP32, $48.00 \pm 1.47\%$ 8-bit, and $43.42 \pm 1.30\%$ 4-bit PTQ). While PolyMerge represents the strongest empirical baseline in this Vision Transformer setup, our proposed wave-theoretic frameworks achieve highly competitive and robust performance. 

\textbf{Analysis of the PolyMerge Empirical Gap.} The empirical superiority of PolyMerge over PhaseMerge (a gap of approximately $5.17\%$ to $6.00\%$ absolute accuracy) can be explained by several architectural and theoretical factors:
\begin{enumerate}
    \item \emph{Macroscopic Layer-Wise Coupling vs. Decoupled Phase Shifts:} PolyMerge constrains its merging coefficients across the model's depth using a low-degree polynomial of the layer index. This matches the macroscopic hierarchical nature of deep transformer networks, where adjacent attention and MLP layers exhibit strong functional dependencies and smooth variations in representation. In contrast, PhaseMerge and U-PhaseMerge learn individual, decoupled phase-shift parameters $\tilde{\phi}_k^l$ for each layer. While this is highly parameter-efficient (192 parameters for $r=1$), the lack of smooth cross-layer coupling can lead to local, high-frequency optimizer drift across layers, slightly degrading global multi-task coordination.
    \item \emph{Spatial Coordinate Dependency on Dense Weights:} As proven in Section 3.4, a uniform phase rotation in Fourier space is mathematically dual to a spatial rotation involving the task vector and its directional Hilbert transform. Because dense layers lack physical coordinates, taking the Hilbert transform of their weight matrices relies on the arbitrary initial neuron indexing. Although this shared basis is locked, mixing a weight update with its coordinate-dependent Hilbert transform can introduce fine-grained high-frequency noise that slightly hampers unquantized performance compared to direct real-space linear scaling.
\end{enumerate}

\textbf{The Hybridization Path: PolyPhaseMerge.} To combine the best of both worlds, we propose \textbf{PolyPhaseMerge} as a highly promising future direction. Instead of optimizing decoupled layer-wise phase angles $\phi_k^l$, we can parameterize the phase shifts $\phi_k(l)$ across layers as a continuous polynomial of the layer index $l$:
\begin{equation}
    \phi_k(l) = a_k + b_k \cdot l + c_k \cdot l^2
\end{equation}
By enforcing a smooth quadratic polynomial constraint on the frequency-domain phase angles across the network's depth, PolyPhaseMerge would integrate PolyMerge's macroscopic depth-wise coordination with PhaseMerge's microscopic phase synchronization, potentially matching or exceeding PolyMerge's absolute performance while retaining PhaseMerge's outstanding low-bit quantization resilience.

Notably, our primary proposed method, \textbf{Uniform PhaseMerge (U-PhaseMerge, $r=1$)}, achieves $42.83 \pm 1.76\%$ in FP32 and $42.33 \pm 1.76\%$ in 8-bit PTQ, significantly outperforming both the static Uniform TA baseline ($38.25 \pm 1.34\%$ and $37.75 \pm 1.43\%$) and the adaptive real-space AdaMerging baseline ($42.00 \pm 0.89\%$ and $41.67 \pm 1.45\%$). This demonstrates that continuous phase rotation in the Fourier domain is exceptionally competitive with, and can outperform, traditional real-space layer-wise optimizations while utilizing a highly structured, low-dimensional phase representation.

Crucially, the static \textbf{FREE-Merging} baseline performs very poorly, achieving only $27.17 \pm 1.96\%$ in FP32 and $27.17 \pm 2.18\%$ in 8-bit PTQ. The massive $+15.66\%$ absolute performance gain of U-PhaseMerge over FREE-Merging provides highly convincing empirical proof that \emph{adaptive} continuous phase synchronization is vital, and that static frequency filtering collapses because it cannot adjust to the alignment dynamics of conflicting experts. By rotating complex-valued phases in Fourier space, PhaseMerge shifts the optimization from sharp coordinate-fitting to continuous wave coherence, showing that frequency-domain phase rotation is a viable and conceptually rich alternative for model merging.

\subsection{Overfitting-Optimizer Paradox (Sample Complexity Sweep)}

To validate the generalizability of PhaseMerge in data-scarce regimes, we sweep the calibration sample size $M \in \{4, 16, 32\}$ and report the resulting multi-task average accuracy under target 8-bit weight quantization.

\begin{table}[ht]
\caption{The Overfitting-Optimizer Paradox: Multi-Task test accuracy (\%) under 8-bit quantization as a function of the calibration stream size $M$.}
\label{tab:sample_complexity}
\vskip 0.15in
\begin{center}
\begin{small}
\begin{sc}
\resizebox{\columnwidth}{!}{%
\begin{tabular}{lccc}
\toprule
Method & $M=4$ & $M=16$ & $M=32$ \\
\midrule
Uniform TA & 37.75 $\pm$ 1.43\% & 37.75 $\pm$ 1.43\% & 37.75 $\pm$ 1.43\% \\
FREE-Merging & 27.17 $\pm$ 2.18\% & 27.17 $\pm$ 2.18\% & 27.17 $\pm$ 2.18\% \\
AdaMerging & 40.75 $\pm$ 1.24\% & 41.67 $\pm$ 1.45\% & 42.50 $\pm$ 1.59\% \\
PolyMerge & \textbf{47.25 $\pm$ 1.87\%} & \textbf{48.00 $\pm$ 1.47\%} & \textbf{47.83 $\pm$ 1.03\%} \\
U-PhaseMerge ($r=1$, Ours) & 42.42 $\pm$ 1.64\% & 42.33 $\pm$ 1.76\% & 40.67 $\pm$ 3.65\% \\
PhaseMerge ($r=2$, Ours) & 40.58 $\pm$ 1.74\% & 40.83 $\pm$ 1.18\% & 42.00 $\pm$ 1.34\% \\
\bottomrule
\end{tabular}%
}
\end{sc}
\end{small}
\end{center}
\vskip -0.1in
\end{table}

\begin{figure}[ht]
    \centering
    \includegraphics[width=0.48\textwidth]{../results/fig2_overfitting_paradox.png}
    \caption{The Overfitting-Optimizer Paradox: Multi-Task test accuracy under 8-bit weight quantization plotted against calibration stream size $M$ with standard deviation error bars over 3 independent random seeds. PolyMerge retains the highest overall performance, while our proposed PhaseMerge variants demonstrate strong calibration capabilities under scarce data.}
    \label{fig:overfitting_paradox}
\end{figure}

The empirical sweep shown in Table~\ref{tab:sample_complexity} and plotted in Figure~\ref{fig:overfitting_paradox} provides striking validation of the Overfitting-Optimizer Paradox and highlights the significant stabilizing effect of our newly-introduced symmetry-preserving frequency mask. At extreme data scarcity ($M=4$), the unconstrained linear optimization of AdaMerging yields $40.75 \pm 1.24\%$, while PolyMerge registers $47.25 \pm 1.87\%$. Our primary proposed formulation U-PhaseMerge ($r=1$), equipped with our mathematical conjugate symmetry mask, achieves a highly competitive $42.42 \pm 1.64\%$, successfully outperforming AdaMerging and closing the gap to PolyMerge. This proves that enforcing zero-phase rotations on DC and Nyquist frequencies physically stabilizes optimization under extreme data scarcity, preventing the Overfitting-Optimizer Paradox.

As the calibration stream size increases to $M=16$, PolyMerge achieves $48.00 \pm 1.47\%$, while U-PhaseMerge ($r=1$) reaches $42.33 \pm 1.76\%$. At $M=32$, PhaseMerge ($r=2$) stabilizes at $42.00 \pm 1.34\%$. This highlights that restricting phase adjustments to preserve conjugate symmetry acts as a physical coordinate regularizer, dramatically smoothing the optimization space compared to unconstrained layer-wise optimization baselines.

\subsection{Target Quantization Schema Shift}

Finally, we evaluate the generalizability of optimized configurations under target deployment schema shifts. Parameters are optimized using an 8-bit quantization operator during calibration but evaluated under different target schemas (4-bit PTQ, 8-bit PTQ, and FP32).

\begin{table*}[ht]
\caption{Target Schema Shift: Multi-Task average accuracy (\%) across different deployment schemas after calibrating under an 8-bit schema.}
\label{tab:schema_shift}
\vskip 0.15in
\begin{center}
\begin{small}
\begin{sc}
\resizebox{\textwidth}{!}{%
\begin{tabular}{lcccccc}
\toprule
Deployment Schema & Uniform TA & FREE-Merging & AdaMerging & PolyMerge & U-PhaseMerge ($r=1$) & PhaseMerge ($r=2$) \\
\midrule
4-bit PTQ & 33.17 $\pm$ 1.39\% & 24.33 $\pm$ 2.37\% & 37.50 $\pm$ 1.22\% & \textbf{43.42 $\pm$ 1.30\%} & 37.42 $\pm$ 1.94\% & 36.92 $\pm$ 0.92\% \\
8-bit PTQ & 37.75 $\pm$ 1.43\% & 27.17 $\pm$ 2.18\% & 41.67 $\pm$ 1.45\% & \textbf{48.00 $\pm$ 1.47\%} & 42.33 $\pm$ 1.76\% & 40.83 $\pm$ 1.18\% \\
FP32 (Unquantized) & 38.25 $\pm$ 1.34\% & 27.17 $\pm$ 1.96\% & 42.00 $\pm$ 0.89\% & \textbf{48.00 $\pm$ 1.62\%} & 42.83 $\pm$ 1.76\% & 40.75 $\pm$ 1.43\% \\
\bottomrule
\end{tabular}%
}
\end{sc}
\end{small}
\end{center}
\vskip -0.1in
\end{table*}

\begin{figure}[ht]
    \centering
    \includegraphics[width=0.48\textwidth]{../results/fig3_schema_shift.png}
    \caption{Target Schema Shift: Multi-Task test accuracy vs.\ target deployment schema bit-depth with standard deviation error bars over 3 independent random seeds, demonstrating that both PolyMerge and PhaseMerge generalize effectively across different target schemas.}
    \label{fig:schema_shift}
\end{figure}

The results in Table~\ref{tab:schema_shift} and Figure~\ref{fig:schema_shift} demonstrate that PolyMerge maintains the highest absolute performance across FP32 and 8-bit target deployment configurations (reaching $48.00 \pm 1.62\%$ and $48.00 \pm 1.47\%$, respectively). Both AdaMerging and our proposed PhaseMerge variants exhibit strong generalizability to target schema shift. When shifting from the 8-bit calibration schema to a 4-bit deployment schema, AdaMerging's performance drops by $4.17\%$, while U-PhaseMerge's performance decreases by only $4.91\%$ and PhaseMerge's performance decreases by only $3.91\%$. This confirms that PhaseMerge produces smooth phase transformations in the frequency domain that are highly robust and do not cause severe quantization noise amplification under shift, validating the generalizability of continuous wave-theoretic optimization under quantization constraints.

\subsection{Ablation Study: Spectral Grid Dimension}
\label{sec:ablation_grid}

To investigate whether the 2D spatial frequency smoothing grid ($r = 2$) provides beneficial spatial regularizations compared to a simpler, global uniform phase-shift, we conduct an ablation study evaluating a $1 \times 1$ grid size ($r = 1$). In this $r = 1$ regime, a single scalar phase-shift is learned per layer and task, which is uniformly broadcast to all frequency components. This represents a highly structured, low-dimensional phase-shift configuration, reducing the optimization parameter footprint from $768$ variables to an extremely compact $192$ variables (matching AdaMerging's footprint but operating in the complex frequency domain).

We calibrate both configurations under the standard 8-bit Post-Training Quantization regime with $M=16$ calibration samples and report multi-task average accuracies across the FP32, 8-bit PTQ, and 4-bit PTQ deployment schemas in Table~\ref{tab:ablation_grid}.

\begin{table}[ht]
\caption{Ablation Study: Multi-Task average accuracy (\%) of PhaseMerge under grid sizes $r=1$ (192-D uniform phase-shift) and $r=2$ (768-D bilinear frequency smoothing grid) optimized at $M=16$.}
\label{tab:ablation_grid}
\vskip 0.15in
\begin{center}
\begin{small}
\begin{sc}
\resizebox{\columnwidth}{!}{%
\begin{tabular}{lccc}
\toprule
Grid Configuration & FP32 & 8-bit PTQ & 4-bit PTQ \\
\midrule
U-PhaseMerge ($r=1$, Ours) & \textbf{42.83 $\pm$ 1.76\%} & \textbf{42.33 $\pm$ 1.76\%} & \textbf{37.42 $\pm$ 1.94\%} \\
PhaseMerge ($r=2$) & 40.75 $\pm$ 1.43\% & 40.83 $\pm$ 1.18\% & 36.92 $\pm$ 0.92\% \\
\bottomrule
\end{tabular}%
}
\end{sc}
\end{small}
\end{center}
\vskip -0.1in
\end{table}

The ablation results in Table~\ref{tab:ablation_grid} reveal distinct, highly informative trade-offs between uniform and spatially continuous phase rotations. Our primary proposed \textbf{U-PhaseMerge ($r=1$)} consistently achieves higher accuracy across all configurations compared to the larger $2\times 2$ grid ($r=2$). Specifically, the $r=1$ grid improves accuracy by $+2.08\%$ in FP32, $+1.50\%$ in 8-bit PTQ, and $+0.50\%$ in 4-bit PTQ. 

This is an extremely clean, coherent finding: it demonstrates that learning a single uniform phase-rotation per layer and task, by being extremely parameter-efficient (192 parameters), acts as a \textbf{much stronger, more robust regularizer} than the 2D spatially-continuous upsampled grid ($r=2$). Because neuron indices in dense layers are unordered, adjacent indices have no spatial correlation, making the upsampled 2D spatial frequency smoothing of $r=2$ somewhat coordinate-dependent. In contrast, the $r=1$ uniform phase shift restricts the optimization to a single shared global degree of freedom per layer. While the 2D FFT itself remains coordinate-dependent due to its complex exponential basis construction, locking the coordinate ordering to the shared pre-trained initialization backbone $W_{\text{pre}}$ allows U-PhaseMerge to act as a stable and robust matrix-basis regularizer, avoiding high-frequency overfitting without introducing arbitrary spatial-frequency smoothing assumptions on the neuron coordinates.

\subsection{Limitations and Future Directions}
While PhaseMerge presents a highly creative conceptual leap, we acknowledge several empirical and theoretical limitations:
\begin{itemize}
    \item \textbf{Sandbox Evaluation Scale:} Due to the rapid prototyping constraints of our sandbox environment, the downstream experts were fine-tuned on 500 samples, and evaluation test sets were subsampled to 100 samples per task. While this is sufficient to prove conceptual validity and execute fast CPU-based runs, scaling PhaseMerge to standard benchmarks evaluated over full test sets remains an important future milestone.
    \item \textbf{Evaluation on Massive Backbones:} Evaluating PhaseMerge using state-of-the-art, fully converged multi-billion parameter Large Language Models (LLMs) or large-scale Vision Transformers will provide a more comprehensive view of its scaling and generalizability potential in production settings.
    \item \textbf{Optimization Instability \& Phase Regularization:} The slight performance variations highlight the non-convex, unstable nature of learning continuous complex variables using test-time entropy objectives. To address this, we incorporated an $L_2$ phase decay penalty ($\mathcal{L}_{\text{reg}} = \gamma \sum \|\tilde{\phi}_k^l\|_2^2$ with $\gamma = 10^{-4}$) directly into our objective function. This penalty constrains the phase-rotation parameters to prevent excessive drift from the initial Task Arithmetic coordinates ($\phi \approx 0$). Empirically, this soft regularization significantly stabilized the optimization dynamics and reduced high-variance outcomes under larger calibration streams ($M=32$), demonstrating the importance of proximity-constrained wave superposition.
\end{itemize}
